\section{Conclusion and Future Work}
Our main finding is that compressing a model in the frequency domain, called band-limiting, gracefully degrades the predictive accuracy as a function of the compression rate.
In this study, we also develop principled schemes to reduce the resource consumption of neural network training.
Neural networks are heavily over-parametrized and modern compression techniques exploit this redundancy.
Reducing this footprint during training is more challenging than during inference due to the sensitivity of gradient-based optimization to noise.

While implementing an efficient band-limited convolutional layer is not trivial, one has to exploit conjugate symmetry, cache locality, and fast complex arithmetic, no additional modification to the architecture or training procedure is needed. Band-limited training provides a continuous knob to trade-off resource utilization vs. predictive performance.
But beyond computational performance, frequency restriction serves as a strong prior.
If we know that our data has a biased frequency spectra or that the functions learned by the model should be smooth, band-limited training provides an efficient way to enforce those constraints.

There are several exciting avenues for future work. Trading off latency/memory for accuracy is a key challenge in streaming applications of CNNs, such as in video processing. Smooth tradeoffs allow an application to tune a model for its own Quality of Service requirements. One can also imagine a similar analysis with a cosine basis using a Discrete Cosine Transform rather than an FFT. There is some reason to believe that results will be similar as this has been applied to input compression~\cite{dctUber2018NIPS} (as opposed to layer-wise compression in our work). Finally, we are interested in out-of-core neural network applications where intermediate results cannot fit in main-memory. Compression will be a key part for such applications. We believe that compression can make neural network architectures with larger filter sizes more practical to study.

We are also interested in the applications of Band-limited training to learned control and reinforcement learning problems.
Control systems are often characterized by the impulse response of their frequency domains.
We believe that a similar strategy to that presented in this paper can be used for more efficient system identification or reinforcement algorithms.